\label{cap:conclusao}

Este trabalho investigou, até o momento, o estado da arte sobre a operacionalização de princípios éticos em sistemas de Inteligência Artificial no contexto da Engenharia de Software, com foco em mecanismos práticos capazes de apoiar desenvolvedores na verificação da conformidade ética ao longo do ciclo de vida do software. Para isso, foi conduzido um Mapeamento Sistemático da Literatura (MSL), abrangendo o período de 2020 a 2025, conforme as diretrizes metodológicas de Petersen et al.~\cite{Petersen2015}. Esse mapeamento constituiu a primeira fase da abordagem metodológica adotada, fundamentando-se na necessidade de consolidar os princípios éticos mais prevalentes na literatura e de identificar lacunas técnicas que orientassem o desenvolvimento dos artefatos propostos nesta dissertação.

Os resultados do MSL mostraram que a literatura sobre ética aplicada à IA concentra-se principalmente em três princípios centrais: \textit{Fairness}, \textit{Accountability} e \textit{Transparency}, identificados em 82,8\%, 55,2\% e 50,0\% dos estudos analisados, respectivamente. Esses princípios não apenas estruturam o núcleo conceitual das investigações recentes, como também são aqueles para os quais há maior maturidade em termos de métodos, \textit{frameworks} e ferramentas de apoio. Ainda assim, a análise detalhada revelou que a operacionalização ética permanece fragmentada, com predominância de abordagens voltadas para as fases de \textit{Requirements}, \textit{Design} e \textit{Implementation}, enquanto a fase de \textit{Testing} apresenta baixa cobertura (37,9\%), mesmo sendo crítica para a detecção de falhas éticas antes da implantação.

De modo geral, os achados do mapeamento evidenciam que, embora exista um corpo crescente de métodos, guias e ferramentas, ainda se carece de soluções integradas e automatizadas que permitam verificar, de forma sistemática e escalável, violações éticas em modelos de IA. Entre as principais lacunas identificadas, destacam-se: (i) ausência de mecanismos de rastreabilidade que conectem decisões éticas ao longo do SDLC; (ii) escassez de ferramentas capazes de operar em ecossistemas reais de desenvolvimento (IDEs, pipelines de CI/CD, plataformas de MLOps); (iii) sub-representação de abordagens aplicadas à fase de testes; e (iv) falta de evidências empíricas robustas sobre a adoção industrial de práticas éticas. Tais lacunas justificam a proposta desta dissertação e orientam a construção das \textit{fuzzing engines} dedicadas aos princípios de \textit{Transparency}, \textit{Fairness} e \textit{Accountability}.

Com base no diagnóstico fornecido pelo MSL, a pesquisa avançou para sua segunda fase, dedicada à identificação e seleção de riscos éticos passíveis de serem operacionalizados por meio de técnicas de \textit{fuzzing}. Essa etapa permitiu estabelecer uma taxonomia de riscos testáveis: dos 11 riscos identificados, 6 são plenamente implementáveis via \textit{fuzzing} automatizado e 5 não são implementáveis por esta abordagem, requerendo métodos complementares de avaliação como auditorias qualitativas, estudos com usuários e análise por especialistas em ciências sociais. Essa taxonomia serve de base para a especificação técnica, implementação modular e posterior avaliação empírica das engines propostas. A especificação técnica adotou o formato padronizado estabelecido no projeto de pesquisa coordenado, no qual esta dissertação se insere, definindo para cada risco implementável: técnica de \textit{fuzzing} adequada, oráculo formal, critérios de aprovação/reprovação e métricas de avaliação.

A análise integrada de métodos de auditoria ética, ferramentas de \textit{fairness}, recursos de explicabilidade e abordagens de governança revelou oportunidades concretas para a aplicação de \textit{fuzzing} como técnica inovadora de testes éticos em IA, especialmente por sua capacidade de explorar automaticamente entradas diversificadas, detectar comportamentos inesperados e estimular condições extremas que podem revelar falhas éticas ocultas.

Portanto, a conclusão parcial deste trabalho indica que há uma lacuna clara e relevante na literatura e na prática da Engenharia de Software quanto à disponibilidade de ferramentas automatizadas para a avaliação sistemática de princípios éticos em IA. A partir desse diagnóstico, as etapas subsequentes da pesquisa concentrar-se-ão em: (i) implementar \textit{fuzzing engines} dedicadas aos princípios de \textit{Transparency}, \textit{Fairness} e \textit{Accountability}; (ii) integrar as engines à arquitetura do sistema de \textit{fuzzing} ético do projeto coordenado; e (iii) conduzir uma avaliação empírica com os modelos ChatGPT (OpenAI), Claude (Anthropic), Gemini (Google) e DeepSeek, a fim de verificar a eficácia das engines e gerar evidências sobre a prevalência de falhas éticas em sistemas de IA contemporâneos.

Os resultados esperados incluem a disponibilização de ferramentas modulares e reutilizáveis, integráveis à arquitetura do projeto coordenado, um conjunto expressivo de casos de teste ético, métricas de conformidade compatíveis com as métricas agregadas do sistema (\textit{failure rate} por risco, modelo e princípio) e diretrizes para adoção prática. Espera-se que esses artefatos contribuam tanto para avanços científicos quanto para a melhoria dos processos de desenvolvimento e auditoria ética de sistemas de IA na indústria.